\begin{frame}{Structure de l'en-tête TEI}
\begin{block}{}
    \textbf{teiHeader}: en-tête contenant les métadonnées décrivant le document.
\end{block}
    \begin{block}{}
        \textbf{4 composants principaux:}
    \begin{itemize}
            \item \texttt{<fileDesc>}~: description d'un document \textcolor{red}{(obligatoire!)};
            \item \texttt{<encodingDesc>}~: description de l'encodage. Relation entre le document électronique et la ou les source(s) dont il est issu;
            \item \texttt{<profileDesc>}~: description non-bibliographique (langue du texte, contexte de production etc...);
            \item \texttt{<revisionDesc>}~: historique des modifications du fichier.
            
    \end{itemize}
    \end{block}
    
\end{frame}

\begin{frame}{Description d'un fichier avec \texttt{<fileDesc>}}
\textbf{Trois parties obligatoires}~:
\begin{itemize}
    \item \texttt{<titleStmt>}~: le titre de l'oeuvre;
    \item \texttt{<publicationStmt>}~: information sur la publication de l'oeuvre (date, lieu...);
    \item \texttt{<sourceDesc>}~: information sur la source dont est tirée l'oeuvre (description du manuscrit par exemple).
\end{itemize}
    
\end{frame}

\begin{frame}[fragile]{Description d'un fichier avec \texttt{<fileDesc>}}
 \begin{lstlisting}[numbers=none]
    <TEI xmlns="http://www.tei-c.org/ns/1.0">
        <teiHeader>
            <fileDesc>
                <titleStmt> Titre </titleStmt>
                <publicationStmt> Date et lieu de la publication</publicationStmt>
                <sourceDesc> Description de la source </sourceDesc>
            </fileDesc>
        </teiHeader>
        <text>
        </text>
    </TEI>
    \end{lstlisting}
    
\end{frame}

\begin{frame}{Description des choix d'encodage avec \texttt{<encodingDesc>}}
    \begin{itemize}
        \item \texttt{<projectDesc>}: description du projet;
        \item \texttt{<editorialDecl>}: précisions sur les pratiques éditoriales employées;
        \item \texttt{<variantEncoding>}: utilisé pour l'encodage d'apparat critique.
    \end{itemize}
\end{frame}

\begin{frame}{Description non-bibliographique avec \texttt{<profileDesc>}}
\begin{itemize}
    \item \texttt{<creation>}~: origine du texte;
    \item \texttt{<langUsage>}~: langue d'écriture du texte;
    \item \texttt{<textDesc>}~: description de texte;
    \item \texttt{<particDesc>}~: description des participants réels ou fictifs;
    \item \texttt{<settingDesc>}~: description du contexte.
\end{itemize}
    
\end{frame}

\begin{frame}{Exercice: en-tête TEI du \textit{Bourgeois Gentilhomme}}
Ecrire le \texttt{<teiHeader>} avec les éléments suivants et en s'appuyant sur les guidelines~:
\begin{itemize}
    \item \texttt{<fileDesc>} (et ses trois parties obligatoires)
            \item \texttt{<encodingDesc>}
            \item \texttt{<profileDesc>}
            \item \texttt{<revisionDesc>}
\end{itemize}
\end{frame}
